\chapter*{Resumen}
\addcontentsline{toc}{chapter}{Resumen}

SavePoint nace de aplicar al videojuego el patrón de catalogación personal de plataformas como Goodreads o Letterboxd, con un inventario de propiedad física y digital que esos productos no necesitan resolver; sobre esa base, el trabajo convierte la plataforma en un banco de pruebas reproducible para sistemas de recomendación. Se ha diseñado e implementado el producto con catálogo, biblioteca, copias físicas y digitales, valoraciones, listas, comentarios, perfiles, privacidad, relaciones sociales y recomendaciones. La aplicación y la evaluación reproducible de los recomendadores son las dos partes complementarias del trabajo: la primera recoge y gobierna los datos que la segunda necesita.

El sistema compara dieciséis algoritmos de referencia: baselines, variantes basadas en contenido, filtrado colaborativo por vecinos e híbridos con reordenación por diversidad. Se describen la representación de obras y perfiles, la ponderación de etiquetas, la similitud por facetas, la calidad externa, la popularidad, la evidencia negativa, el arranque en frío y la combinación híbrida. El protocolo offline fija corpus, snapshots, semillas, candidatas comunes, exclusiones y métricas para $K\in\{5,10,20\}$.

La evidencia publicada v15 utiliza 400 usuarios sintéticos activos y 79 usuarios evaluables, con una sola ejecución del test. Los resultados permiten comparar internamente algoritmos bajo el corpus y protocolo congelados, pero no representan usuarios reales ni justifican generalizaciones automáticas. La memoria distingue de forma explícita esa publicación del contrato de implementación vigente, que evolucionó posteriormente.

Se ha aplicado una metodología de ingeniería asistida por agentes, gobernada mediante el método Get Stuff Done y organizada conceptualmente con Obsidian. Las herramientas se emplearon bajo dirección y revisión humana, manteniéndose la decisión final y la responsabilidad sobre el análisis, el diseño, la implementación y la validación en la dirección del proyecto. La trazabilidad entre requisitos, código, pruebas, decisiones y artefactos versionados constituye una aportación metodológica del trabajo.

\vspace{.5cm}
\textbf{Palabras clave:} videojuegos, sistemas de recomendación, evaluación offline, reproducibilidad, filtrado colaborativo, recomendación basada en contenido, procedencia de datos.
